\documentclass[12pt]{article}

% ── Packages ──────────────────────────────────────────────────────────────────
\usepackage[margin=1.25in]{geometry}
\usepackage{setspace}

\usepackage{amsmath,amssymb}
\usepackage{mathtools}
\usepackage[hidelinks]{hyperref}
\usepackage[numbers,sort&compress]{natbib}
\usepackage{booktabs}
\usepackage{enumitem}
\usepackage{titlesec}
\usepackage{abstract}
\usepackage{xcolor}
\usepackage{parskip}

% ── Typography ─────────────────────────────────────────────────────────────────
\setstretch{1.15}
\setlength{\parskip}{0.6em}
\setlength{\parindent}{0pt}

\titleformat{\section}{\normalfont\large\bfseries}{{\thesection}}{1em}{}
\titleformat{\subsection}{\normalfont\normalsize\bfseries}{{\thesubsection}}{1em}{}
\titlespacing*{\section}{0pt}{1.4em}{0.5em}
\titlespacing*{\subsection}{0pt}{1em}{0.3em}

\renewcommand{\abstractnamefont}{\normalfont\bfseries}
\renewcommand{\abstracttextfont}{\normalfont\small}

% ── Math shorthand ─────────────────────────────────────────────────────────────
\newcommand{\sS}{\mathcal{S}}
\newcommand{\sT}{\mathcal{T}}
\newcommand{\sG}{\mathcal{G}}
\newcommand{\reach}{\mathrm{Reach}}

% ──────────────────────────────────────────────────────────────────────────────
\begin{document}

\begin{center}
  {\LARGE\bfseries From Preservation to Reachability:}\\[0.4em]
  {\large\bfseries Semantic Continuity in an Age of Drift}\\[1.2em]
  {\normalsize Flyxion}\\[0.2em]
  {\small\textit{Independent Researcher} \quad
   \texttt{github.com/standardgalactic}}\\[1.0em]
  {\small April 17, 2026}\\[0.2em]
  {\small\textit{revised June 4, 2026}}
\end{center}

\begin{abstract}
\noindent
Contemporary information systems are typically evaluated in terms of
preservation: archives store documents, platforms retain content, and
institutions maintain records.
This paper argues that preservation is neither the correct invariant nor
the correct design objective for semantic continuity.
What matters is whether semantic states remain \emph{reachable} from future
states through available transformations, and whether such reachability
continues to support processes of reconstruction and reconstitution.
We formalize the distinction between preservation and reachability, introduce
\emph{ecphory} as the activation condition required for semantic reconstruction,
and analyze a series of failure modes in which semantic states persist while
their accessibility collapses.
Extending the framework to language, memory, archives, note systems, and
cultural transmission, we argue that semantic continuity depends on the
maintenance of continuation paths rather than the storage of isolated states.
A field-theoretic interpretation is proposed in terms of scalar persistence,
vector transport, and accessibility fields, leading to a broader conclusion:
what survives across time is not primarily information or representation, but
the possibility that fragments may encounter compatible contexts and become part
of larger reconstructed wholes.
\end{abstract}

\newpage

% ── 1. Introduction ────────────────────────────────────────────────────────────
\section{Introduction}

A library that burns is understood as a catastrophe.
A library whose catalog is deleted, whose call numbers are reassigned, and whose
shelves are rearranged without documentation is understood as---
at worst---a management failure.
Yet in both cases the result is the same: a reader arriving with a reference
cannot find the book.
The second library has \emph{preserved} its collection while destroying the
conditions under which the collection is useful.

This contrast captures the central problem of the present paper.
Modern digital and institutional systems are organized around preservation:
the persistence of files, pages, references, and audiences.
But preservation, in the sense of continued existence, is neither necessary nor
sufficient for what we actually require from memory systems, which is
\emph{reachability}: the ability of agents in future states to arrive at past
semantic content through available transformations.

The distinction matters because preservation and reachability can fail
independently.
A state can be preserved but rendered unreachable (the archive that no search
engine indexes).
A path can exist but lead nowhere (the link that resolves to a renamed page
with overwritten content).
And a system can maintain all its original components while gradually
eliminating the connections among them.

The argument proceeds in stages.
Section~2 introduces the formal distinction between preservation and reachability
and extends it to the concept of reconstitution.
Section~3 develops reconstitution as the purpose of reachability and introduces
the principle that a fragment encounters a compatible context and becomes
something larger than itself.
Section~4 introduces ecphory as a third condition, distinct from both
preservation and reachability, required for reconstruction to succeed in
practice.
Section~5 presents a taxonomy of failure modes---each a distinct way
preservation and reachability can come apart---together with the neighborhood
principle as a constructive counterpart.
Section~6 extends the analysis to natural language systems, where the same
structure operates at the scale of generations rather than server cycles.
Section~7 sketches the constructive implication: what a system must do if it
aims to preserve reachability rather than merely storage.
Section~8 develops a field-theoretic interpretation of the three constructive
conditions, introducing scalar persistence, vector transport, and accessibility
fields.
Section~9 and~10 extend the analysis to forgetting, the structural role
of incompleteness as a continuation operator, and the dynamics of semantic
momentum.
Section~11 develops the argument's final implication: that the correct
invariant for semantic continuity is not preservation or reachability
but continuation.

A single formulation crystallizes the distinction before the formalism: a system
can \emph{remember perfectly} while being \emph{unable to recall}.
Memory, in this sense, is the existence of a state in the system.
Recall is the existence of a path from here to there.
The two are independent, and conflating them is the source of most design
failures in knowledge infrastructure.

% ── 2. Preservation and Reachability ──────────────────────────────────────────
\section{Preservation and Reachability}

Let a \emph{semantic system} consist of a set of states $\sS$ and a set of
transformations $\sT$ that agents can apply to move between states.
These transformations include any operation that constitutes legitimate access or
navigation in the system: following a hyperlink, querying a search index, looking
up a citation, using a translation, recalling from a prior interaction.
The pair $(\sS, \sT)$ induces a directed graph $\sG$ in which states are nodes
and transformations are edges.

\textbf{Preservation} is the property that a state $s \in \sS$ continues to
exist over time---that is, that the node remains in the graph.

\textbf{Reachability} is the property that a state $s$ remains in the forward
closure of the present state $s_0$ under available transformations---that is,
that there exists a directed path from $s_0$ to $s$ in $\sG$.

The central claim of this paper is:

\begin{quote}
\textit{A system can preserve elements of $\sS$ while destroying the path
structure of $\sG$, such that preserved states become unreachable from the
present.}
\end{quote}

This claim is not a logical paradox but a structural observation about how
systems are actually built.
Nodes and edges are maintained by different mechanisms---storage systems versus
indexing and linking systems---and they degrade through different processes.
Preservation is addressed by backup policies, archival standards, and data
redundancy.
Reachability is addressed by navigation infrastructure: search, citation,
cross-referencing, naming consistency.
Because these are separate engineering concerns, they can fail separately.

Floridi's account of semantic information treats information as meaningful and
truthful data~\cite{floridi2011}, but the question of \emph{access} to
information is largely treated as prior to or outside the information-theoretic
frame.
Shannon's channel model~\cite{shannon1948} similarly brackets semantic content
in favor of transmission fidelity.
What neither framework captures directly is the temporal and navigational
dimension: whether a receiver in a future state can reconstruct the path to a
past message through the transformations available at that future time.
Reachability, as we use the term here, names this gap.

\subsection{Reachability Is Agent-Relative}

Throughout this paper, reachability has been treated as a property of a
semantic graph together with an agent class capable of traversing that
graph.
This dependence deserves explicit emphasis before the analysis proceeds.

A path that exists for one agent may not exist for another.
The trained archivist, the domain expert, the native speaker, and the
long-term participant in a community possess transformations unavailable
to outsiders.
They inhabit richer effective graphs.
Consequently, accessibility is not properly defined as a function of
states alone; it must be indexed by an agent class $A$:
\[
  \mathcal{A}(s, A, t) \;=\; \log \lvert \mathcal{P}(s, A, t) \rvert.
\]
The same archive, language, or note system therefore contains multiple
accessibility fields simultaneously, one for each class of agent capable
of traversing it.

This observation clarifies why preservation projects often appear
successful from the perspective of specialists while failing from the
perspective of the general public.
The preserved states remain reachable for the small population that already
knows how to navigate the system, even as their accessibility collapses
for everyone else.

Semantic continuity is therefore never absolute.
It is always continuity \emph{for some population of traversing agents}.
A system that narrows the agent class capable of traversal---through
increased credential requirements, specialist vocabulary, or the
disappearance of communities with relevant navigational knowledge---reduces
semantic continuity even when every node and every edge is formally intact.
The degradation is invisible to standard preservation metrics because it
occurs not in the stored states but in the population of agents for whom
those states remain live.

% ── 3. Reachability and Reconstitution ────────────────────────────────────────
\section{Reachability and Reconstitution}

The formal framework of Section~2 defines reachability as the existence
of a path from a present state to a past state.
This definition is necessary but not yet sufficient to explain why
reachability matters.
The answer requires a further step.

Reachability matters because it enables \emph{reconstitution}: the
process by which a partial present state, brought into contact with a
preserved trace through an available transformation, generates a larger
semantic structure than either component contained alone.

A preserved document is not valuable because it exists.
It is valuable because it can participate in reconstructing a meaning
that a present agent could not have generated independently.
A citation is valuable not because it points to a node but because
following the citation enables a reader to reconstruct an argument,
a method, or a context that reshapes what the reader can subsequently
understand.
A language is valuable not because it encodes a lexicon but because
it enables speakers to reconstruct the conceptual structures that prior
speakers developed over generations.

This allows us to sharpen the paper's central claim.
The failure of preservation systems is not merely that they allow states
to become unreachable.
It is that they allow \emph{reconstitution events to become impossible}:
future agents arrive with questions that prior states could answer, but
the paths that would carry those agents to the relevant states are gone.

We may state the underlying principle as follows:

\begin{quote}
\textit{A fragment encounters a compatible context and becomes something
larger than itself.}
\end{quote}

Reachability is the condition under which this encounter can occur.
A preserved but unreachable state is a fragment that can never encounter
a compatible context again.
It persists as inscription while ceasing to participate in the ongoing
generation of meaning.

This framing also clarifies the relationship between the six cases.
Each failure mode is not merely a failure of access; it is a failure
of reconstitution.
Namespace laundering makes the reconstruction of prior community meaning
impossible for agents who approach through the current name.
Encyclopedic consolidation makes the reconstruction of prior framings
impossible for agents who arrive through general-purpose search.
Access restriction makes reconstruction impossible for agents who lack
credentials.
Link rot makes it impossible for anyone.
Archival isolation makes it impossible for all but specialists who
already possess the navigational knowledge to reach the archive.

The dense local system (Section~5.6), conversely, succeeds not because
it stores well but because it continuously enables reconstitution: present
members encounter past states through living practice, and the encounter
regenerates understanding that no document alone could transmit.

% ── 4. Ecphory and Semantic Recall ────────────────────────────────────────────
\section{Ecphory and Semantic Recall}

The framework developed so far identifies three conditions for semantic
continuity: node persistence, edge maintenance, and path coherence.
But these conditions are necessary, not sufficient.
Even when all three are satisfied---a state exists, a path exists, and the
path is coherent---reconstruction may still fail.

This gap is captured by the concept of \emph{ecphory}: the event in which
a present cue resonates with a preserved trace strongly enough to
reconstruct a larger semantic state.

Ecphory names the \emph{activation event}, distinct from both preservation
and reachability.
A state may be preserved.
A path to it may exist.
Yet if the present state does not carry sufficient energy, context, or
structural compatibility to resonate with the trace, the reconstruction
does not occur.
The path exists in principle but is not traversed in practice.

This gives a three-layer framework:

\begin{enumerate}[leftmargin=2em,itemsep=0.2em]
  \item \textbf{Preservation}: the state continues to exist ($s \in \sS$).
  \item \textbf{Reachability}: a path from the present to the state exists
        ($s \in \reach(s_0)$).
  \item \textbf{Ecphory}: the present state activates the preserved trace,
        enabling reconstruction of a larger semantic whole.
\end{enumerate}

Each condition can fail independently.
A state can be preserved but unreachable (archival isolation).
A state can be reachable but not ecphorically activated---the path is
available, but the agent's current state does not carry sufficient
resonance to complete the reconstruction.
A cue can be present and a path can be followed, but if the surrounding
context lacks the structural compatibility to reconstruct the original
encoding conditions, the traversal terminates without reconstitution.

The ecphoric failure mode is distinct from the infrastructural failure modes
catalogued in Section~5.
It is not a failure of storage or infrastructure; it is a failure
of resonance.
The agent must not only be able to reach the state; they must arrive in
a condition capable of activating it.

This has practical implications for system design.
A retrieval system optimized for reachability delivers agents to the
relevant node.
A system optimized for ecphoric activation additionally ensures that
agents arrive with the contextual preparation required for reconstruction.
The difference is visible in the contrast between a search result
(reachability) and a well-designed curriculum (ecphoric preparation):
the curriculum does not merely provide access to content; it sequences
prior encounters so that each state the learner has traversed increases
their capacity to activate the next.

The ecphoric framework also explains a pattern in the failure modes.
Encyclopedic consolidation (Section~5.2) is particularly damaging not
only because it removes intermediate nodes but because it removes the
intermediate states that provided the ecphoric preparation required to
activate the absorbed content.
A reader who arrives at a consolidated article without having traversed
the prior framing may be able to read the words but unable to reconstruct
the understanding---because the encoding context that made the original
formulation meaningful has been stripped away along with the intermediate
nodes.

\subsection{Ecphoric Failure and the Tip-of-the-Tongue State}

The three-layer framework---preservation, reachability, ecphory---predicts
a class of failure that is distinct from all the infrastructure failure
modes catalogued in Section~5.

Consider a state $s_m$ that is preserved ($s_m \in \sS$), reachable from
the present ($s_m \in \reach(s_0)$), and yet not reconstructed.
The path exists; the agent can in principle follow it; but the present
state does not carry sufficient resonance with the target trace to
complete the ecphoric activation.

This failure mode has a precise phenomenological signature.
It appears subjectively as \emph{partial accessibility}: the agent knows
that the target state exists, senses its proximity, and may even retrieve
peripheral features---associated concepts, partial phonological or visual
form, contextual fragments---without being able to reconstruct the state
itself.

This is the tip-of-the-tongue phenomenon, and under the present framework
it is not a retrieval failure.
It is an ecphoric threshold failure.

Preservation remains intact: the trace is stored.
Reachability may remain intact: a path from the present state to the
target exists.
What fails is activation: the resonance between the current cue and the
stored trace is insufficient to complete the reconstruction.

The practical consequence is that introducing additional energy into
the retrieval process---a secondary cue, a change of context, an
interval of rest that allows the retrieval wave to propagate through
intermediate states without active interference---can supply the
additional resonance required for ecphoric completion.
This explains why tip-of-the-tongue states frequently resolve
spontaneously, often after attention has been redirected elsewhere:
the retrieval process continues subthreshold, traversing intermediate
states until the accumulated resonance crosses the activation boundary.

The tip-of-the-tongue case makes concrete why reachability alone is
insufficient as an account of semantic continuity.
A system could be perfectly navigable---every state reachable from
every other through short paths---and still fail to reconstruct a
target state if the agent's current configuration lacks the structural
compatibility required for ecphoric activation.
Design for semantic continuity must therefore attend not only to path
structure but to the conditions under which agents arrive at states
prepared to reconstruct rather than merely to traverse.

% ── 5. A Taxonomy of Discontinuities ──────────────────────────────────────────
\section{A Taxonomy of Discontinuities}

The following cases are not chosen for comprehensiveness but for
structural distinctness.
The first five each represent a different way in which the reachability graph
$\sG$ can be damaged independently of the state set $\sS$.
The sixth inverts the pattern, showing that reachability can be sustained
without formal storage infrastructure at all.
Together the six constitute what we will call the \emph{minimal generating set}
of cases for understanding semantic continuity.

\subsection{Namespace Laundering: Identity Drift with Audience Retention}

On large social platforms, named communities---groups, pages, channels---
accumulate audiences over time.
The audience constitutes the social infrastructure through which content
is distributed and meaning is contextually stabilized.
When a community is renamed or its stated purpose is changed, the audience
is typically retained while the semantic identity of the community is
overwritten.

This operation, which we call \emph{namespace laundering}, severs the mapping
between a name and its prior referent while conserving the distribution
network.
Audience members who joined under one semantic contract now participate in
one that has been silently revised.
New members encounter the community only in its present form, with no record
that a prior form existed.

The failure here is not loss of a node but corruption of the edge label.
The path from a present state to the community's prior meaning passes through
a name that now points elsewhere.
The prior meaning is not preserved; it is overwritten at the level of the
identifier.

\subsection{Encyclopedic Consolidation: Path Compression with Loss}

Reference systems such as Wikipedia stabilize meaning through consolidation:
articles are merged, disambiguated, or renamed as editors converge on
preferred framings.
This process improves navigability in one direction---from the present toward
canonical content---while eliminating intermediate nodes that represented
prior framings, minority perspectives, or earlier consensus states.

When an article is merged into another, the prior article's existence is
recorded only in edit histories, which are not indexed by external search
and require specific navigational knowledge to access.
A user arriving from a search query encounters the current consensus without
any indication that a prior structure existed.
The prior structure is preserved in the sense that it can be reconstructed
by someone who already knows to look, but it is not reachable from the
present by available general-purpose transformations.

This is path compression with loss: the graph is simplified, but the
simplification is irreversible from the outside.

\subsection{Access Restriction: Gated Edges}

Academic publishing systems preserve a large proportion of the scholarly
record.
Digital Object Identifiers~\cite{paskin2010} provide persistent references
that survive server changes and link decay.
Articles exist, are indexed, and can be cited.
But without institutional subscription credentials, the edge from a reference
to its referent is blocked at the point of traversal.

The node is present in $\sG$.
The edge exists.
But the edge is \emph{gated}: traversal is restricted to agents who satisfy
an external condition unrelated to the semantic content of the document.
From the perspective of an agent without credentials, the state is
preserved but unreachable.
The citation points to an object that cannot be accessed through available
transformations.

\subsection{Link Rot: Edge Deletion}

Web infrastructure links are edges, and edges degrade.
Studies of link persistence in scholarly and journalistic contexts consistently
find substantial decay rates over periods of years~\cite{klein2014}.
Domains lapse, servers are decommissioned, and URL structures are reorganized
without redirection.
The result is a reference that resolves to nothing: a citation without an
object, a pointer with no target.

This is the most straightforward case---literal edge deletion.
The node may have existed; it may no longer.
What is certain is that no path from the present state leads to it.
The failure is not at the level of preservation policy but at the level of
infrastructure maintenance: the transformations that once existed no longer
do.

\subsection{Archival Isolation: Node Preservation with Path Loss}

Archival projects represent the most serious attempt to address
preservation---the deliberate maintenance of digital states against decay.
But preservation of a node is not sufficient if no paths lead to it.

The September 11 Digital Archive, maintained by the Roy Rosenzweig Center for
History and New Media and held by the Library of Congress, preserves a
substantial record of firsthand accounts and digital artifacts from 2001.
It is not indexed by major search engines in ways that surface it for
relevant queries.
Users who do not already know of its existence typically discover it only
through secondary references---citations in academic work, mentions in
archival literature, or the Wayback Machine's own indexing.

The archive is preserved.
It is not reachable from the present by the transformations available to a
general-purpose user.
The path exists---in principle---but requires specific navigational knowledge
that is itself subject to decay.

This case reveals a deeper point: reachability is not a binary property of
a system but a function of the navigational resources available to a
particular class of agent at a particular time.
More precisely, reachability is a property not of the graph $\sG$ alone,
but of the pair $(\sG, A)$, where $A$ denotes the capabilities of a given
agent---the transformations that agent can actually perform.
Two agents inhabiting the same system may therefore inhabit different
effective graphs: the credentialed researcher and the uncredentialed reader
face different subgraphs of $\sG$, as do the trained archivist and the
general-purpose search user, or the fluent speaker and the monolingual
outsider.
An archive reachable by specialists is not reachable in the sense that
matters for broad semantic continuity.

\subsection{Dense Local Systems: Reachability Without Archival Memory}

The five cases above are all failure modes.
It is instructive to examine a contrasting case---a system in which
reachability is effectively guaranteed not by archival infrastructure
but by path density.

Consider a small, stable community with high identity continuity: a
neighborhood, a craft guild, a religious congregation, a family that has
occupied the same region for generations.
Such communities typically have weak formal archival infrastructure.
They do not maintain indexed records of past decisions or digitized
repositories of prior knowledge.
Yet semantic continuity is often extremely robust.

The reason is structural: the reachability graph $\sG$ is locally dense.
Every present state is connected to past states through multiple redundant
paths---overlapping memory among members, repeated narrative, ritual
re-enactment, shared material culture, and the direct presence of
individuals who participated in prior events.
If one path degrades, others remain.
The system does not need formal preservation because it continuously
traverses its own history.

This case makes the theoretical point most cleanly: storage is not the
invariant.
A system with minimal storage and maximal path density may sustain semantic
continuity more robustly than a system with comprehensive archives and
degraded navigation infrastructure.
The dense local case is not a nostalgic ideal; it is a structural limit
that clarifies what the failure modes are actually failing to reproduce.
What large-scale information systems lose when they scale is not primarily
storage capacity---it is exactly this local redundancy of paths.

\subsection{The Neighborhood Principle}

The dense local system case suggests a general principle that applies
across scales: reachability is a function of neighborhood structure,
not merely of node existence.

Let $d(s)$ denote the \emph{local path density} of state $s$: the number
of independent paths connecting $s$ to other states within a bounded
traversal radius.
Then the effective reachability of $s$ is better characterized by $d(s)$
than by the binary fact of $s$'s existence.

\[
  \reach(s) \;\approx\; f\bigl(d(s)\bigr)
\]

A state embedded in a dense neighborhood---surrounded by many other states
that agents traverse in the course of ordinary activity---is reached
repeatedly as a byproduct of other searches.
A state in a sparse neighborhood is reached only by agents who specifically
seek it, and only if those agents already possess the navigational knowledge
to approach it.

This has direct implications for note-taking, archival design, and
knowledge management.
A note that is placed beside many other frequently-visited notes acquires
reachability through the traffic of the neighborhood.
A note filed in isolation---even if correctly categorized---acquires
reachability only through direct search, which requires the searcher to
already know what they are looking for.

The practical design principle follows:
\begin{quote}
\textit{Do not file a note where it belongs. Place it where it can
become relevant.}
\end{quote}

The two criteria diverge whenever the future retrieval context differs
from the classification context.
Taxonomies organize states according to present classification---what
a thing \emph{is} at the moment of filing.
Neighborhood placement organizes states according to anticipated future
traversal---\emph{how} and \emph{from where} a thing will be encountered.
``Where it belongs'' is determined by the taxonomy at the moment of
filing.
``Where it can become relevant'' is determined by the predicted traffic
patterns of the neighborhoods it might inhabit.
When these differ---as they routinely do in any evolving intellectual
project---the neighborhood placement is the more durable choice.

The same principle operates at larger scales.
A language community's semantic heritage is reachable not because
it is stored but because it is embedded in the neighborhoods that
living speakers traverse daily---in idioms, rituals, canonical
texts, and shared narrative.
When the community disperses and those neighborhoods dissolve, the
heritage becomes unreachable not because the records have been
destroyed but because the path density has collapsed.

% ── 6. Language as a Long-Timescale Instance ──────────────────────────────────
\section{Language as a Long-Timescale Instance}

The five failure modes described above operate at the timescale of years and
decades.
The same structure is visible in natural language systems, but the timescale
is generations, and the mechanisms are social and political rather than
technical.

Language is not simply a medium of communication; it is a navigation system
for a vast inherited semantic field.
A speaker of a language inherits access to a particular subgraph of human
meaning---the texts, oral traditions, conceptual structures, and
interpersonal knowledge encoded in that language's history.
A language shift---whether voluntary or imposed---is not merely a change in
expressive medium.
It is a change in the available transformations: the paths that remain
navigable after the shift are determined by what has been translated,
preserved, or actively maintained across the discontinuity.

The displacement of minority languages by dominant ones---French displacing
regional languages in metropolitan France, English reorganizing the
linguistic landscape of colonial territories, Spanish restructuring prior
Mesoamerican semantic systems---is precisely a restructuring of the
reachability graph.
Prior states (oral traditions, untranslated texts, contextual knowledge
embedded in grammatical structure) are not destroyed in the moment of
displacement.
They become unreachable by the paths available to subsequent generations,
who inherit the new navigational infrastructure.

The persistence of Arabic as a living navigational system across fourteen
centuries of geographical dispersal and political transformation is an
instructive counterexample.
Arabic maintains its reachability through continuous traversal: liturgical
practice, Quranic recitation, classical literary education, and an
unbroken tradition of commentary and cross-referencing.
The classical texts are not merely preserved; they are \emph{continuously
reached}, repeatedly, by a large population of agents who maintain the
transformations that make them accessible.
This is active path maintenance rather than archival storage.

The contrast between the two cases---a language displaced without maintained
paths versus a language kept reachable through continuous practice---
illuminates the constructive requirement.
It is not enough to record a language; the navigational infrastructure must
also be sustained.

The same logic applies to institutional language policy.
Quebec's \emph{Charte de la langue fran\c{c}aise} (1977) is typically analyzed
as a policy of linguistic protection or identity assertion.
From the present perspective, it is better understood as a constraint on
reachability: it enforces, through legal means, the maintenance of edges in
the francophone semantic graph that would otherwise be eliminated by
the dominant pressure of English-language commercial and administrative
infrastructure.
The policy does not restore already-lost paths; it acts as a boundary
condition that prevents certain transitions from occurring, maintaining
connectivity to prior semantic states that would otherwise become
unreachable within a generation.

This is not an analogy to the technical case.
It is the same mechanism operating through different implementation media.
The reachability graph of a language is damaged by the same structural
operations---edge deletion, path compression, audience redirection---that
damage the reachability graphs of digital knowledge systems.
The timescales differ; the underlying structure does not.

% ── 7. The Constructive Requirement ───────────────────────────────────────────
\section{The Constructive Requirement}

The analysis above yields a requirement, not merely a diagnosis.
If the goal of a memory system is semantic continuity---the ability of future
agents to reach past meaning through available transformations---then the
design question shifts from ``how do we store information?'' to ``how do we
ensure that future states can reach past states through valid paths?''

This is a stronger condition than preservation.
It requires attending to three things simultaneously:

\begin{enumerate}[leftmargin=2em,itemsep=0.2em]
  \item \textbf{Node persistence}: states must continue to exist.
  \item \textbf{Edge maintenance}: the transformations that connect states must
        remain available and traversable.
  \item \textbf{Path coherence}: the composition of available transformations
        must remain capable of connecting present states to past ones.
\end{enumerate}

Current systems are organized primarily around the first condition.
Archival standards, redundant storage, and persistent identifiers address node
persistence.
The second and third conditions are largely left to market incentives
(search engine indexing), institutional inertia (citation practices), and
individual effort (translation, commentary, cross-referencing).

The asymmetry is not incidental.
Node persistence is a technical property that can be specified, measured, and
contractually guaranteed.
Edge maintenance and path coherence are relational properties that require
ongoing coordination among the agents who constitute the system.
They cannot be archived; they must be continuously performed.

This points toward a broader framework in which meaning is treated not as a
stored object but as a property of a system under transformation---something
that must remain \emph{invariant under the transformations that the system
makes available to its agents}.
Constraint-based accounts of meaning, in which semantic content is defined
by what it rules out rather than what it encodes, capture this naturally:
to preserve meaning is to preserve the constraints that make certain
transitions valid and others not.
When those constraints are lost---when the paths are severed---the content
may persist as a physical inscription while losing its operative semantic
character.

This connection is developed in Section~8, which introduces a field-theoretic
decomposition of the reachability condition and reinterprets the failure modes
geometrically under that framework.
What the present analysis establishes is the conceptual prior step: that
reachability, not preservation, is the invariant that matters for semantic
continuity, and that modern information systems are systematically organized
to maintain the wrong thing.

\subsection{Continuation Requires Friction}

Modern information systems frequently treat friction as an obstacle to be
eliminated.
Search engines minimize traversal effort, recommendation systems compress
path length, and interfaces attempt to remove every intermediate step
between question and answer.

Yet the complete removal of friction produces a paradox.
Traversal does not merely recover semantic states; it partially constructs
them.
The intermediate encounters, contextual cues, and neighboring structures
experienced during navigation contribute to the conditions required for
ecphoric activation.
A path of zero length may maximize retrieval efficiency while minimizing
the contextual preparation that makes reconstruction possible.

The dense local systems discussed in Section~5.6 illustrate the opposite
extreme.
Their semantic continuity derives not from frictionless access but from
repeated traversal through overlapping pathways.
Meaning is maintained because agents continuously move through the graph
rather than arriving directly at isolated destinations.
The intermediate steps are not wasted effort; they are the mechanism by
which each traversal reinforces the accessibility of surrounding states
and prepares the agent for ecphoric activation at the target.

The objective of a continuation-oriented system is therefore not the
elimination of friction but the cultivation of \emph{productive friction}:
enough traversal to maintain contextual coupling and neighborhood density
without introducing so much resistance that paths become inaccessible.

This suggests a design criterion beyond the three conditions of
node persistence, edge maintenance, and path coherence.
A fourth condition might be called \emph{traversal adequacy}: the system
must route agents through a sufficient density of intermediate states
that each arrival is prepared to reconstruct rather than merely to
reach.
Curricula, annotation systems, guided tours, and contextual recommendation
all serve this function---not by shortening paths but by ensuring that
the paths agents traverse build the contextual state required for
meaningful reconstruction at the destination.

% ── 8. Toward a Field-Theoretic Account ───────────────────────────────────────
\section{Toward a Field-Theoretic Account}

The three constructive conditions identified in Section~7 have a natural
interpretation that extends beyond graph theory.
Node persistence, edge maintenance, and path coherence are not merely
engineering requirements; they correspond to distinct physical quantities
in any system where semantic content propagates through a medium rather
than residing at fixed addresses.

We sketch this correspondence here not as a full formal development but
as an orientation toward the richer treatment the reachability framework
points toward.

\subsection{The Scalar--Vector--Entropy Decomposition}

Consider a \emph{continuation field} over the semantic manifold, decomposed
into three components:

\begin{itemize}[leftmargin=2em,itemsep=0.3em]
  \item A \textbf{scalar persistence field} $\Phi(x,t)$, recording the
        local density of states that continue to exist at position $x$
        and time $t$.
        High $\Phi$ means the neighborhood is well-populated with
        preserved nodes.
        This is the field-theoretic form of node persistence.

  \item A \textbf{vector transport field} $\mathbf{v}(x,t)$, recording
        the direction and magnitude of available transformations at each
        point---the flow structure that carries agents from one state
        to another.
        Degraded $\mathbf{v}$ corresponds to gated edges, link rot, and
        path compression: the transformations nominally exist but do not
        propagate effectively.
        This is the field-theoretic form of edge maintenance.

  \item An \textbf{accessibility field} $\mathcal{A}(x,t)$, recording
        the volume of admissible future trajectories emanating from
        each state.
        Formally, if $\mathcal{P}(x)$ denotes the set of valid continuation
        paths from state $x$ through available transformations, the
        accessibility is:
        \[
          \mathcal{A}(x,t) \;=\; \log \lvert \mathcal{P}(x,t) \rvert.
        \]
        High $\mathcal{A}$ means the state participates in many possible
        futures---it can be reached from many directions and leads to many
        continuations.
        Low $\mathcal{A}$ means the state is a near-isolated node: preserved
        but functionally inert.
        This is the field-theoretic form of path coherence.
\end{itemize}

Under this decomposition, the central claim of the paper becomes precise.
A system can maintain high $\Phi$ (storage is intact) while allowing
$\mathbf{v}$ and $\mathcal{A}$ to collapse.
The archive that no search engine indexes has high $\Phi$ and
near-zero $\mathcal{A}$.
The gated journal article has high $\Phi$, an existing edge in $\mathbf{v}$,
and an accessibility that varies discontinuously with agent credentials.
The dense local community has modest $\Phi$ and very high $\mathcal{A}$:
redundant paths maintain reachability even without formal storage.

The accessibility field $\mathcal{A}$ is the quantity the paper has been
tracking throughout, under the name reachability.
The contribution of the field-theoretic framing is to make explicit that
$\mathcal{A}$ is not a binary property of nodes but a continuous function
over the semantic manifold---one that can degrade smoothly, concentrate
in specialist subgraphs, or be maintained by active traversal even in
the absence of archival infrastructure.

\subsection{Geometric Reinterpretation of the Failure Modes}

The six cases of Section~5 admit a uniform geometric description under
this framework.

\textbf{Namespace laundering} is a discontinuous transformation of the
vector field: the edge labels are reassigned so that flows that formerly
reached a prior state are redirected to a new one.
$\Phi$ is unchanged; $\mathbf{v}$ is redirected; $\mathcal{A}$ of the
prior state collapses as incoming paths are severed.

\textbf{Encyclopedic consolidation} is a contraction: intermediate nodes
are removed and paths are compressed.
The resulting graph has lower local path multiplicity.
$\mathcal{A}$ decreases for the absorbed states, which now require
specialist navigational knowledge---a narrowing of the effective agent
class $A$---to reach.

\textbf{Access restriction} introduces a potential barrier into the
vector field: the edge exists in $\sG$ but is traversable only by
agents whose capability set $A$ satisfies an external credential condition.
For uncredentialed agents, the effective transport field $\mathbf{v}$
contains a gap at the point of restriction.

\textbf{Link rot} is literal deletion of transport channels.
Edges disappear from $\mathbf{v}$; no flow can cross the gap;
$\mathcal{A}$ of the downstream states falls to zero for agents
approaching from the deleted direction.

\textbf{Archival isolation} maximizes $\Phi$ while allowing $\mathbf{v}$
and $\mathcal{A}$ to degrade through neglect of indexing and
cross-referencing infrastructure.
The state exists; the paths to it do not.

\textbf{Dense local systems} exhibit the inverse: modest $\Phi$
compensated by high local path density.
Multiple redundant edges in $\mathbf{v}$ maintain high $\mathcal{A}$
even when individual paths degrade.
This is the structural reason such systems sustain semantic continuity
robustly without formal archives: the accessibility field is maintained
by continuous traversal rather than by storage policy.

\subsection{Accessibility as a Continuous Quantity}

The graph-theoretic framework of Section~2 treats reachability as a
binary property: a state is either reachable from the present or it is
not.
This discretization is useful for stating the central claim precisely,
but it misrepresents the actual structure of semantic systems.

Real semantic systems exhibit \emph{graded accessibility}.
A state is not simply reachable or unreachable; it is reachable by more
or fewer agents, through more or fewer paths, with greater or lesser
effort, at higher or lower probability of successful ecphoric activation.
The graph is a discrete approximation of an underlying continuous
accessibility field.

Formally, the directed graph $\sG$ may be understood as a discretization
of a smooth manifold equipped with a flow structure.
The nodes are samples from a continuous state space; the edges are
discrete renderings of what is in reality a continuous transport field
$\mathbf{v}(x,t)$.
The binary reachability predicate $s \in \reach(s_0)$ is a threshold
applied to the continuous accessibility functional
$\mathcal{A}(x,t) = \log|\mathcal{P}(x,t)|$.

This continuity is empirically visible in each domain the paper has
considered.
Archives do not become inaccessible at a single moment; accessibility
decays as link rot accumulates, indexing degrades, and the agent
community that knows how to navigate the archive shrinks.
Languages do not become unreachable overnight; the path density
thins over generations as speaker communities disperse and
transmission practices weaken.
Memories do not disappear; they become progressively harder to
activate as the cue-trace resonance falls below ecphoric threshold.
Notes do not vanish; they become functionally inert as the
neighborhoods surrounding them become less frequently traversed.

In each case, the useful question is not ``is this state reachable?''
but ``what is the current value of $\mathcal{A}$ at this state, for
this class of agent, under current field conditions?''
The binary question asks whether a path exists.
The continuous question asks how many paths exist, how robust they are
to perturbation, how large an agent class can traverse them, and how
much energy a traversal requires.

This continuous framing also clarifies the relationship between the
failure modes and cognitive states.
The tip-of-the-tongue state---the subjective experience of partial
accessibility, of knowing that a memory exists without being able to
retrieve it---is precisely what high-$\Phi$, low-$\mathcal{A}$ feels
like from the inside.
The trace is preserved.
The paths to it have degraded below the threshold of effortless
traversal.
The state is reachable in principle but requires more energy than the
present cognitive state can supply.

The viscosity of the retrieval medium, the density of the local
neighborhood, and the resonance between current cue and target trace
all determine whether a given state lies above or below the ecphoric
threshold for a given agent at a given moment.
These are continuous quantities, not binary ones, and the framework
is impoverished if they are reduced to a simple present/absent
distinction.

\subsection{Persistence as Continuation}

The deepest implication of the field-theoretic framing is a restatement
of the paper's central claim in terms that generalize across domains.

A system organizes around \emph{preservation} when it treats $\Phi$ as
the invariant to be maintained.
A system organizes around \emph{reachability} when it treats $\mathcal{A}$
as the invariant.

But $\mathcal{A}$ is not a property of nodes.
It is a property of the dynamics of the system---of how states participate
in ongoing processes of traversal, reconstruction, and continuation.
A state with high $\mathcal{A}$ is not merely reachable in principle;
it is actively embedded in a web of possible futures.
It remains coupled to the ongoing evolution of the manifold.

This suggests that the correct invariant is neither preservation nor
reachability in isolation, but what might be called \emph{continuation}:
the property that a semantic state remains embedded in live trajectories
rather than existing as an isolated inscription.

Continuation fails when $\mathbf{v}$ degrades (the transport channels
decay), when $\mathcal{A}$ collapses (the state ceases to participate
in possible futures), or when the agent class $A$ capable of performing
the relevant traversals shrinks below the threshold required to maintain
active path density.
Preservation addresses none of these failure modes.
Active path maintenance, cross-referencing, translation, and institutional
constraint on edge deletion address all three.

The connection to constraint-based accounts of meaning noted in Section~7
becomes clearer under this framing.
Meaning is preserved when the constraints that govern valid transitions
are maintained---when the rules that make certain paths traversable and
others not continue to operate across time and across the agent population
that constitutes the system.
When those constraints erode, content persists as physical inscription
while its operative semantic character---its participation in live
continuation---disappears.

The formal elaboration of this principle, including coupled field equations
for $(\Phi, \mathbf{v}, \mathcal{A})$ under various degradation and
maintenance regimes, lies beyond the scope of the present paper.
What the analysis here establishes is the conceptual prerequisite:
that any adequate account of semantic continuity must track not only
whether states exist, but whether they remain embedded in ongoing
processes of reaching and being reached.

\subsection{Continuation as a Dynamical Quantity}

The accessibility field $\mathcal{A}(x,t)$ captures the current volume of
admissible continuations available from a state.
Yet accessibility alone is insufficient to characterize the long-term
behavior of semantic systems.

A state may possess high accessibility while contributing little to future
continuation.
Conversely, a state with modest present accessibility may generate large
numbers of future traversals through its influence on subsequent states.

This suggests that continuation should be treated as a dynamical quantity
rather than a static one.

Let $C(x,t)$ denote the \emph{continuation potential} of a state: the
expected future production of accessibility generated by traversals
through that state.
Where $\mathcal{A}$ measures the accessibility currently available,
$C$ measures the accessibility likely to be produced.

The distinction is analogous to that between stored energy and power.
Two states may possess identical accessibility while differing radically
in their capacity to generate future accessibility.
A forgotten archive may possess substantial latent accessibility but
generate little future traversal.
A foundational theorem, a canonical story, or a persistent unresolved
question may continually regenerate accessibility by routing future
agents through itself and into the surrounding neighborhood.

Continuation therefore depends not only on the present geometry of the
semantic manifold but on the manifold's capacity to reproduce that
geometry through ongoing traversal.

Under this interpretation, the semantic momentum introduced in
Section~10.1 becomes a special case of continuation dynamics.
Momentum measures observed traversal persistence at a given time.
Continuation potential $C(x,t)$ measures the underlying capacity to
generate such persistence into the future.

The transition from preservation to reachability and from reachability
to continuation may therefore be understood as a transition from static
state variables to dynamical ones.
Preservation asks whether a state exists.
Reachability asks whether a path exists.
Continuation asks whether the system actively regenerates the conditions
required for future paths to exist.

This framing points toward the coupled field equations deferred earlier.
A complete dynamical theory would specify how $\Phi$, $\mathbf{v}$,
$\mathcal{A}$, and $C$ evolve together---how traversal builds or erodes
transport channels, how accessibility collapses or self-reinforces, and
whether critical transitions or phase changes occur as the continuation
potential of a semantic system falls below threshold.
That theory lies beyond the present paper, but the variables it requires
are now defined.

% ── 9. Forgetting as Path Degradation ─────────────────────────────────────────
\section{Forgetting as Path Degradation}

The infrastructural and linguistic failure modes catalogued in Sections~5
and~6 have a cognitive counterpart.
Forgetting, in the standard model, is the loss of stored content: the
memory trace degrades, the information disappears, the record is gone.
The reachability framework suggests a different account.

Forgetting is not primarily the loss of stored content.
Forgetting is the loss of reconstructive pathways.

The evidence for this distinction is familiar but rarely formalized.
People regularly report memories that were ``gone'' returning intact
after years of inaccessibility, triggered by a smell, a phrase, a
location, or an emotional state that reinstates the original encoding
context.
The memory was not absent; it was unreachable.
The retrieval cue did not create the memory; it restored a path that
had degraded below traversal threshold.

This distinction has a precise formulation in the reachability framework.
Let $s_m$ denote a stored memory state.
Traditional forgetting is the event $s_m \notin \sS$: the node disappears.
Reachability-based forgetting is the event that $s_m \in \sS$ but
$s_m \notin \reach(s_0)$: the node persists while the paths from any
present state $s_0$ to $s_m$ become impassable.

The mechanisms of cognitive path degradation are structurally identical
to those in digital infrastructure.
\emph{Cue extinction} is the cognitive analog of link rot: the retrieval
cues that formerly connected present states to $s_m$ are no longer
available (the smell is no longer encountered, the social context has
dissolved, the language of encoding is no longer spoken).
\emph{Interference} is the analog of namespace laundering: competing
associations redirect the retrieval flow, so that paths formerly reaching
$s_m$ now terminate at a different attractor.
\emph{Contextual shift} is the analog of encyclopedic consolidation:
the encoding context has been subsumed into a new frame, and the original
intermediate states are no longer accessible by general-purpose recall.

The practical consequence is that the appropriate response to forgetting
is not re-encoding but path restoration.
Re-encoding adds a new node; it does not restore the paths that made
the original node accessible.
Path restoration reinstates the cues, contexts, and associative
neighborhoods that carry retrieval waves from present states to the
target.
This is why retrieval practice---repeated traversal of the path from
present state to memory---is more effective for long-term retention
than additional study: practice maintains the edges, not just the nodes.

The cognitive analog of the dense local system (Section~5.6) is the
well-rehearsed skill or deeply familiar domain.
In such systems, semantic content is maintained not through deliberate
storage but through continuous traversal.
The craftsperson who has practiced a skill for decades does not
``remember'' how to do it in the sense of consulting a stored procedure;
they maintain the skill because the relevant paths are continuously
active.
Forgetting occurs not when the trace is lost but when the practice ceases
and the paths begin to degrade.

This account also clarifies the relationship between external memory
systems and internal recall.
A note placed in a navigable neighborhood---beside other notes that
will be encountered in the ordinary course of future work---is not
merely external storage.
It is path maintenance: the note continuously reinstates a retrieval
route that would otherwise degrade.
The failure mode of the forgotten note is not that the information
disappeared; it is that the surrounding navigational context decayed,
and no path from any likely future state leads to the note any longer.

% ── 10. Incomplete Structures as Continuation Operators ───────────────────────
\section{Incomplete Structures as Continuation Operators}

The analysis of forgetting reveals a structural asymmetry that applies
more broadly: the relationship between closure and reachability is not
monotone.
Finished artifacts tend to reduce their own accessibility over time.
Incomplete artifacts tend to preserve it.

This is not a productivity claim or a celebration of disorganization.
It is a structural observation about how accessibility is maintained in
semantic systems.

A fully resolved note, essay, or argument reaches a stable attractor.
It says what it says.
The connections it makes are made; the questions it poses are answered;
the tensions it introduces are discharged.
As a semantic object, it is closed.
A closed object participates in a fixed set of future trajectories---those
defined by its explicit content and the links it currently possesses.

An incomplete structure behaves differently.
An unresolved question, an unfinished argument, a stub entry, a marginal
annotation that breaks off mid-thought: each of these maintains an open
edge in the reachability graph.
The incompletion is not a gap in storage; it is a forward hook---a
directed connection into a future state space that has not yet been
determined.

We may formalize this as follows.
Let $\partial s$ denote the \emph{open boundary} of a semantic state $s$:
the set of transformation slots that remain unfilled, questions that remain
unanswered, or connections that are indicated but not completed.
Then the accessible future state space of $s$ grows with $|\partial s|$.
A state with $|\partial s| = 0$ is closed; its continuation set is
determined entirely by its explicit content and current links.
A state with large $|\partial s|$ remains coupled to a larger space of
possible futures---it can be completed in many ways, and different
completions are activated by different subsequent contexts.

This gives a precise sense in which certain forms of incompleteness
increase long-term reachability.
A question survives longer than an answer because it remains connected
to a larger future state space.
A question that has been answered is closed; the answer either persists
in the reader's accessible graph or it does not.
A question that remains open continues to attract future traversal
from anyone who encounters the same problem.

The same principle explains a pattern in intellectual history.
Foundational texts that remain generative across centuries are typically
not texts that resolve everything.
They are texts that open more problems than they close---that establish
a framework rich enough to generate ongoing questions.
The texts that are merely correct tend toward closure and obsolescence.
The texts that remain productive tend to maintain large $|\partial s|$:
they are dense with incompletion, with indicated but unfollowed paths,
with problems named but not solved.

In practical terms, this suggests that note systems should not optimize
exclusively for closure.
A note that records a resolved insight is valuable; a note that records
an unresolved tension is often more valuable, because the tension
continues to attract future attention and to generate new connections.
The marginalia, the stub, the half-written question, and the unfinished
draft are not failures of organization.
They are externalized continuation operators: structures that maintain
open edges in the reachability graph and keep future trajectories
available.

The optimal note, under this analysis, does three things simultaneously:
it records a semantic trace (node persistence), it situates the trace
within a navigable neighborhood (edge maintenance), and it maintains an
open boundary that keeps the trace coupled to future contexts (path
coherence through continuation).
A finished note satisfies the first two conditions.
A strategically incomplete note satisfies all three.

% ── 10.1 Semantic Momentum ────────────────────────────────────────────────────
\subsection{Semantic Momentum}

Not all states contribute equally to future continuation.
Some states generate unusually large numbers of subsequent traversals.
A frequently cited paper, a foundational text, a central theorem, a
canonical example, or a persistent unanswered question may continue to
attract agents long after its original context has disappeared.

We describe this tendency as \emph{semantic momentum}.

Let $\mathcal{T}(s, t)$ denote the expected future traversal rate through
state $s$ at time $t$.
States with high semantic momentum continuously regenerate their own
accessibility by attracting new paths and reactivating old ones.
Their ongoing traversal maintains high $\mathcal{A}$ even in the absence
of deliberate preservation effort.
States with low semantic momentum depend on external maintenance to remain
reachable; without active curation, their accessibility fields decay.

This distinction clarifies a pattern in intellectual history that
preservation-centered accounts cannot easily explain.
Many carefully preserved artifacts become functionally inert: they exist,
they are indexed, paths to them are maintained, yet they attract no
traversal and participate in no reconstructions.
Meanwhile, other artifacts---sometimes poorly preserved, sometimes
available only in fragmentary form---continue to generate ongoing
interpretation, citation, and reconstruction.

The difference is not archival quality.
It is semantic momentum: the capacity of a state to route future agents
through itself, thereby continuously replenishing the accessibility field
surrounding it.

States with high semantic momentum often share structural features with
the incomplete structures discussed earlier.
They tend to raise questions rather than close them, establish frameworks
that generate further problems, and maintain large open boundaries
$|\partial s|$ that keep them coupled to a wide range of future
investigative trajectories.
The connection is not coincidental: incompleteness and semantic momentum
are related aspects of the same underlying property---the capacity of a
semantic state to remain generative rather than merely persistent.

Continuation therefore depends not only on existing accessibility but on
the capacity of states to \emph{generate future accessibility}.
Some states are merely reachable.
Others actively reproduce the conditions of their own reachability by
continuously attracting the traversal that keeps them live.
A theory of semantic continuity that attends only to preservation and
path structure misses this dynamic: the ongoing generation of
accessibility through traversal is as important as its initial
construction.

% ── 10.2 Continuation Across Scales ──────────────────────────────────────────
\subsection{Continuation Across Scales}

The framework developed in this paper exhibits a property that deserves
explicit statement: it is \emph{substrate-independent}.

The same continuation structure appears at every scale of analysis
considered here.
At the neural and cognitive scale, ecphoric recall requires that retrieval
cues carry sufficient resonance to reactivate preserved traces through the
pathways that connect present states to target attractors.
At the scale of individual knowledge management, notes and artifacts
persist functionally when they are embedded in navigable neighborhoods
that future cognitive states will traverse.
At the institutional scale, archives sustain semantic continuity when
indexing, cross-referencing, and agent communities maintain the transport
infrastructure.
At the linguistic scale, languages remain reachable when communities
continuously traverse the paths that connect present speakers to inherited
semantic content.
At the cultural scale, traditions survive when practices of re-enactment,
commentary, and transmission continuously regenerate the paths between
present participants and prior events.

In each case, the mechanism is the same.
A partial present state---a retrieval cue, a search query, a citation,
a spoken word, a ritual gesture---encounters a preserved trace through
an available transformation, and the encounter enables the reconstruction
of a larger semantic whole that neither component contained independently.

The framework characterizes a \emph{relation} rather than a storage medium.
This is why the same structure is visible across implementations that
differ in almost every other respect: neural dynamics, paper notebooks,
database indices, oral tradition, institutional practice.
The substrate varies; the relational structure does not.

This scale-invariance is not a coincidence.
It reflects the fact that semantic continuity is always a matter of
whether partial structures can encounter compatible contexts---and that
this condition is independent of the physical medium in which the
structures are encoded.
Memory, archives, languages, and institutions are all solutions to the
same problem: how to preserve the conditions under which such encounters
remain possible across time, across agents, and across the inevitable
transformations that any system undergoes.

The principle that states this most concisely:

\begin{quote}
\textit{A fragment encounters a compatible context and becomes something
larger than itself.}
\end{quote}

This is not a metaphor drawn from one domain and applied to others.
It is a structural description of the process by which semantic continuity
is maintained at every scale at which it has been examined here.
The implication is that any system---cognitive, archival, linguistic,
cultural, or institutional---that aims to sustain meaning across time
must preserve not objects but the conditions under which this encounter
remains possible.

% ── 11. From Preservation to Continuation ─────────────────────────────────────
\section{From Preservation to Continuation}

The argument of this paper has proceeded in stages, and it is worth naming
the movement explicitly before closing.

The starting point was the distinction between preservation and reachability:
a system can remember perfectly while being unable to recall.
The failure modes of Section~5 demonstrated that this is not a theoretical
edge case but the normal operating condition of contemporary digital
infrastructure.

The extension to language showed that the same structure operates at the
timescale of generations, through social and political mechanisms rather
than technical ones.
Reachability is not a property of digital systems; it is a property of
any system that separates the storage of semantic states from the
maintenance of paths among them.

The constructive requirement (Section~7) named three conditions---node
persistence, edge maintenance, path coherence---and identified the
asymmetry: systems are organized almost exclusively around the first while
the second and third are left to market incentives and individual effort.

The field-theoretic account (Section~8) translated these conditions into
continuous quantities $(\Phi, \mathbf{v}, \mathcal{A})$, identified the
accessibility functional $\mathcal{A}(x,t) = \log|\mathcal{P}(x,t)|$ as
the quantity the paper has been tracking throughout, and introduced the
continuation potential $C(x,t)$ as the dynamical variable that
distinguishes states that merely persist from those that actively
regenerate the conditions of their own reachability.

The sections on ecphory (Section~4), including the tip-of-the-tongue
analysis (Section~4.1), neighborhood (Section~5.7), forgetting (Section~9),
incompleteness (Section~10), semantic momentum (Section~10.1), and
continuation across scales (Section~10.2)
extended the analysis beyond infrastructure systems to cognitive and cultural
systems, and identified a common mechanism in each: partial structures that
enable or block the reconstitution of larger wholes.

What emerges from this sequence is a shift across three levels of description.

\emph{Preservation} treats semantic content as an object to be stored.
The question is: does the object still exist?

\emph{Reachability} treats semantic content as a node in a navigable graph.
The question is: does a path from here to there still exist?

\emph{Continuation} treats semantic content as a participant in ongoing
processes of traversal and reconstruction.
The question is: does this state remain embedded in live trajectories?

Each level is strictly stronger than the previous.
A state can be preserved without being reachable.
A state can be reachable without being in continuation---if the paths
exist in principle but are never traversed, the state effectively disappears
from the active semantic manifold.

The correct invariant for semantic continuity is therefore not preservation,
not reachability, but continuation: the property that a semantic state
remains coupled to ongoing processes of reconstruction, traversal, and
future possibility.

This reframing has practical implications that follow directly from the
analysis.
Systems designed for preservation optimize storage, redundancy, and
persistence identifiers.
Systems designed for reachability additionally optimize indexing,
cross-referencing, and navigational infrastructure.
Systems designed for continuation do both of those things and additionally
maintain the active agent population whose ongoing traversal keeps
accessibility from collapsing---through institutions, practices, pedagogy,
translation, commentary, and the deliberate cultivation of communities
capable of reaching and reconstructing the semantic states in question.

The deepest lesson of the dense local system case (Section~5.6) is that
continuation does not require formal infrastructure at all.
What it requires is continuous traversal---the ongoing movement of agents
through the paths that connect present states to past ones.
Formal infrastructure is a technology for sustaining traversal at scale
when the community is too large or too distributed for direct path
maintenance.
When that infrastructure fails, it fails by allowing traversal to cease
while giving the appearance that storage alone is sufficient.

The principle that unifies the cases considered here---digital archives,
languages, note systems, ecphoric recall, cultural transmission---may be
stated simply:

\begin{quote}
\textit{A fragment encounters a compatible context and becomes something
larger than itself.}
\end{quote}

Continuation is the condition under which this encounter remains possible.
What must be preserved is not the object, nor even the relation, but the
possibility of the encounter itself: the live embedding of semantic
fragments in the ongoing processes through which meaning is reached,
reconstructed, and transmitted forward in time.

\newpage

% ── References ────────────────────────────────────────────────────────────────
\bibliographystyle{plainnat}

\begin{thebibliography}{9}

\bibitem{shannon1948}
Shannon, C.~E. (1948).
\newblock A mathematical theory of communication.
\newblock \textit{Bell System Technical Journal}, 27(3), 379--423.

\bibitem{floridi2011}
Floridi, L. (2011).
\newblock \textit{The Philosophy of Information}.
\newblock Oxford University Press.

\bibitem{paskin2010}
Paskin, N. (2010).
\newblock Digital object identifier ({DOI}\textregistered) system.
\newblock In \textit{Encyclopedia of Library and Information Sciences},
  3rd~ed. Taylor \& Francis.

\bibitem{klein2014}
Klein, M., Van de Sompel, H., Sanderson, R., Shankar, H., Balakireva, L.,
  Liu, K., \& Moretto, R. (2014).
\newblock Scholarly context not found: One in five articles suffers from
  reference rot.
\newblock \textit{PLOS ONE}, 9(12), e115253.

\end{thebibliography}

\end{document}

